\documentclass[a4paper,11pt]{article}

\usepackage[T1]{fontenc}
\usepackage[margin=2.5cm]{geometry}
\usepackage{amsmath,amssymb}
\usepackage{siunitx}
\usepackage{graphicx}
\usepackage{booktabs}
\usepackage[hidelinks]{hyperref}

\title{Linux Computer Literacy Report}
\author{Hugh Rodgers \\ Student number: 123340696 \\ Email: 123340696@umail.ucc.ie}
\date{\today}

\begin{document}
    \maketitle

    \tableofcontents
    \newpage  


    \section{File Manager - Directories}
        File managers are used to navigate and manage the file system. They allow you to create, move, copy, and delete files and directories efficiently.
        It is very important to keep your files organized and to save data in the appropriate formats to ensure easy access and proper functionality.
        Proper file management helps prevent data loss, improves productivity, and makes it easier to collaborate with others.
        

        \subsection{file types, saving data}
            Certian files are used for certian things, for example:
            \begin{itemize}
                \item \texttt{.txt} for text files
                \item \texttt{.sh} for shell scripts
                \item \texttt{.py} for Python scripts
                \item \texttt{.TEX} for LaTeX documents
                \item \texttt{.md} for Markdown files
                \item \texttt{.csv} for comma-separated values files
                \item \texttt{.bak} for backup files
            \end{itemize}

        creating backups is essential to prevent data loss. Backup files, often with the extension \texttt{.bak}, are copies of important files that can be restored if the original files are accidentally deleted or corrupted.


    \section{Python \& pip}

        \subsection{Installing Python}
        First you need a compiler like VScodium. This can be downloaded and installed from its official website or through your system's package manager and the terminal.
        You can run the command sudo apt install python3 to install Python on your system. Make sure to verify the installation by running python3 --version in the terminal.
        
        \subsection{Python Files}
        .py, .ipynb and other Python-related files are commonly used in Python development. .py files are standard Python scripts, while .ipynb files are Jupyter Notebook files 
        that allow you to create and share documents containing live code, equations, visualizations, and narrative text.


        \subsection{Installing pip}
        To install pip, you can use the following command:
        \begin{verbatim}
        python -m ensurepip --upgrade
        \end{verbatim}
        Note: you may have a different version of python installed just replace \texttt{python} with the appropriate version, e.g., \texttt{python3}.

        \subsection{Installing Packages}
        To install packages using pip, you can use the following command:
        \begin{verbatim}
        pip install "package_name"
        \end{verbatim}
        If a Virtual Environment is not set up this will install the package to the site-packages directory of your Python environment. 



    \section{Venv}

        \subsection{Creating a Virtual Environment}
        A virtual environment is an isolated Python environment that allows you to manage dependencies for different projects separately. You can create a virtual environment 
        using the following command:
        \begin{verbatim}
        python -m venv myenv
        \end{verbatim}

        \subsection{Activating a Virtual Environment}
        To activate the virtual environment, use the following command:
        \begin{verbatim}
        source myenv/bin/activate  # On linux or MacOS
        myenv\Scripts\activate     # On Windows
        \end{verbatim}        
        


    \section{BASH}

        BASH (Bourne Again SHell) is a command-line shell used for interacting with the operating system.
        It allows you to execute commands, run scripts, and automate tasks efficiently.
        the command "- -help" after any command to get more information about its usage. 
        or the command "man" followed by the command name to access the manual. For example:
 
        \subsection{Basic BASH Commands}

            \begin{verbatim}
            # List files and directories
            ls

            # Change directory
            cd /path/to/directory

            # Create a new directory
            mkdir new_directory

            # Remove a file
            rm file_name

            # Remove a directory and its contents
            rm -r directory_name

            # Display the current directory
            pwd

            # stop running processes in the terminal
            Ctrl+C

            # copy and paste changes picking up a shift
            Cntrl+Shift+V  # Paste in terminal
            Cntrl+Shift+C  # Copy in terminal

            # stop the terminal
            exit

            # clear the terminal
            clear

            # run as administrator
            sudo command_here

            # search for a file or directory
            find path -name file_name

            % open vscodium 
            codium . # Open the current directory in VScodium
            codium file_name  # Open a specific file in VScodium
            \end{verbatim}
        
        \subsection{Creating Bash Shortcuts ~/.bashrc file}
        I have added an example ~/.bashrc file in the attached .txt file.
        In linux you can add a ~/.bashrc file which you can use to create custom aliases and functions to make your terminal workflow more efficient.
        This makes setting up a new project very easy, for example when setting this up i was able to put in 3 or 4 command to the
        terminal, which connect this directory to the github repository, created the .TEX files, created the .venv, created the 
        .vscode/settings.json with some preloaded files to hide.



    \section{GIT}

        GIT is a version control system that allows you to track changes in your code and collaborate with others.
        it is usefull for managing code changes and collaborating with others.

        \subsection{Git Commands}

            \begin{verbatim}
                echo "# Example_Project" >> README.md
                git init                        
                # Initialize a new Git repository
                
                git add README.md               
                # Stage the README file for commit

                git commit -m "first commit"    
                # Commit the staged changes with a message

                git branch -M main              
                # Rename the default branch to main

                git remote add origin git@github.com:UserName/Example_Project.git 
                # Add the remote repository
                
                git push -u origin main         
                # Push the changes to the remote repository

                git remote add origin git@github.com:UserName/Example_Project.git 
                # Add the remote repository
                
                git branch -M main              
                # Rename the default branch to main
                git push -u origin main         
                # Push the changes to the remote repository
            \end{verbatim}

            In your  degree you will use git to manage your code, track changes, and collaborate with your peers on various projects




    % \section{Fun with the Terminal}
    %         Here are some fun things you can do with the terminal:
    %         sudo apt install cowsay
    %         sudo apt install fortune
    %         sudo apt install cmatrix
    %         sudo apt install fastfetch or others
    %         sudo apt install sl
    %         sudo apt install wtf            
\end{document}

